\section*{Author source: profunctors and matrix representation (34:4, 34:6--7)}
We write $I$ for the trivial category, which has one object $\star$ and one morphism $\mathrm{id}_\star$. For groups $G$ and $H$, $G\le H$ means that $G$ is a subgroup of $H$.
We write $[0,\infty]$ for the set $\{r\in\mathbb R\mid r\ge0\}\cup\{\infty\}$. The sum $\sum_{j\in J}r_j$ of a possibly infinite family $(r_j)_{j\in J}$ of elements in $[0,\infty]$ is well-defined as the least upper bound of $\{\sum_{k\in K}r_k\mid K\subseteq J,\ K:\text{finite}\}$ (where $\infty+r=\infty$).
A setoid is a set equipped with an equivalence relation. Given a setoid $(X,\approx_X)$, we write $[x]_\approx$ for the equivalence class $\{x'\in X\mid x\approx_Xx'\}$.
A groupoid $A$ is a category in which every morphism is invertible. Then $A(a,a)$ is a group for every object $a\in A$. We use the diagrammatic order for composition: $(\alpha;\beta)\in A(a,c)$ when $\alpha\in A(a,b)$ and $\beta\in A(b,c)$.
For groupoids $A$ and $B$, a profunctor $P:A\pf B$ is a functor $P:A^{\mathrm{op}}\times B\to\mathsf{Set}$. For $x\in P(a,b)$, $\alpha\in A(a',a)$ and $\beta\in B(b,b')$, we write $\alpha\cdot x\cdot\beta$ for $P(\alpha,\beta)(x)$. When $\beta=\mathrm{id}$, we simply write $\alpha\cdot x\cdot\mathrm{id}$ as $\alpha\cdot x$, and similarly for $x\cdot\beta:=\mathrm{id}\cdot x\cdot\beta$.
Profunctors $P:A\pf B$ and $Q:B\pf C$ compose analogously to relations:
\[
(P;Q)(a,c):=\int^{b\in B}P(a,b)\times Q(b,c)
 :=\left(\coprod_{b\in B}P(a,b)\times Q(b,c)\right)/{\sim},
\]
where $\sim$ is the equivalence relation defined by
\[
\operatorname{in}_b(x,y)\sim\operatorname{in}_{b'}(x',y')
\Longleftrightarrow
(x,y)=(x'\cdot\beta,\beta^{-1}\cdot y')\text{ for some }\beta\in B(b',b).
\]
We shall simply write $\operatorname{in}_b(x,y)$ as $(x,y)$. Then $(x\cdot\beta,y)\sim(x,\beta\cdot y)$ for every $x\in P(a,b)$, $\beta\in B(b,b')$ and $y\in Q(b',c)$. The groupoid action is defined by $\alpha\cdot[(x,y)]_\sim\cdot\beta:=[(\alpha\cdot x,y\cdot\beta)]_\sim$.
The weighted relational model with coefficient $[0,\infty]$ is written $\WRel$. Its objects are setoids and its morphisms from $A$ to $B$ are $((A/{\cong_A})\times(B/{\cong_B}))$-matrices $R=(R_{[a],[b]})$ with $R_{[a],[b]}\in[0,\infty]$.
Given $R\in\WRel(A,B)$ and $S\in\WRel(B,C)$, their composite is
\[
(R;S)_{[a],[c]}:=\sum_{[b]\in(B/{\cong_B})}R_{[a],[b]}S_{[b],[c]}
:=\sup\left\{\sum_{[b]\in J}R_{[a],[b]}S_{[b],[c]}
 \,\middle|\,J\subseteq_{\mathrm{fin}}(B/{\cong_B})\right\}.
\]
Here $J\subseteq_{\mathrm{fin}}(B/{\cong_B})$ means a finite subset; $0\cdot\infty:=0$ and $r\cdot\infty:=\infty$ when $r>0$.
\paragraph{Definition 6 (Matrix representation).}
We say a groupoid $A$ is locally finite if $\#A(a,a)<\infty$ for every $a$. A factor of $A$ is an assignment $\Phi^\oplus:\operatorname{obj}(A)\to(0,\infty)$ such that $\Phi^\oplus(a)=\Phi^\oplus(a')$ when $a\cong_A a'$. We define $\Phi^\ominus(a):=\#A(a,a)/\Phi^\oplus(a)$, so $\Phi^\oplus(a)\Phi^\ominus(a)=\#A(a,a)$.
The matrix representation $\norm P$ of a profunctor $P:A\pf B$ between locally finite groupoids $A$ and $B$ with factors is a matrix $\norm P\in\WRel(\norm A,\norm B)$ given by
\begin{equation}
\norm P_{[a],[b]}:=\frac{\#P(a,b)}{\Phi^\ominus(a)\Psi^\oplus(b)},\tag{2}
\end{equation}
where $\norm A$ is the setoid of the objects of $A$.
\section*{Author source: creeds and atomic decomposition (34:9--10)}
\paragraph{Definition 9 (Creed [32]).}
Let $P:I\pf A$. The creed of $P$, written $\mathcal CP$, is a family of subsets $(\mathcal CP)_a\subseteq A(a,a)$ parameterised by objects $a\in A$ defined by
\[
(\mathcal CP)_a:=\{\alpha\in A(a,a)\mid\exists x\in P(a).\ x\cdot\alpha=x\}
=\bigcup_{x\in P(a)}\operatorname{fix}(x),
\]
where $\operatorname{fix}(x)\le A^{\mathrm{op}}(a,a)\times B(b,b)$ for a profunctor $P:A\pf B$ and $x\in P(a,b)$ is the stabiliser subgroup, i.e.,
\[
\operatorname{fix}(x):=\{(\alpha,\beta)\mid\alpha\in A(a,a),\ \beta\in B(b,b),\ \alpha\cdot x\cdot\beta=x\}.
\]
The creed of $Q:A\pf I$ is defined similarly. A family $\mathcal C=(\mathcal C_a)_a$ of subsets $\mathcal C_a\subseteq A(a,a)$ is a creed if $\mathcal C=\mathcal CP$ for some $P:I\pf A$.
Note that we do not define $\mathcal CP$ for general profunctors $P:A\pf B$ (with $A\ne I\ne B$).
\paragraph{Definition 12.} Let $\mathcal C$ and $\mathcal C'$ be creeds on $A$. They are orthogonal, written $\mathcal C\perp\mathcal C'$, if $\mathcal C_a\cap\mathcal C'_a\subseteq\{\mathrm{id}\}$ for every $a$.
Each hom-category $\mathsf{Prof}(A,B)$ has sums defined as follows. For profunctors $P_i:A\pf B$, the profunctor $\coprod_{i\in I}P_i:A\pf B$ is defined by $(\coprod_{i\in I}P_i)(a,b):=\coprod_{i\in I}(P_i(a,b))$ on objects and $(\coprod_{i\in I}P_i)(\alpha,\beta)(\operatorname{inj}(x)):=\operatorname{inj}(P_j(\alpha,\beta)(x))$ on morphisms, where $\operatorname{inj}:P_j(a,b)\to(\coprod_{i\in I}P(a,b))$ is the injection.
A subgroup $G\le A(a,a)^{\mathrm{op}}\times B(b,b)$ induces a profunctor $\widetilde G:A\pf B$ given by
\[
\begin{aligned}
\widetilde G(a',b')&:=\{\alpha\cdot G\cdot\beta\mid\alpha\in A(a',a),\ \beta\in B(b,b')\},\\
\widetilde G(\alpha,\beta)&:=X\mapsto(\alpha\cdot X\cdot\beta),
\end{aligned}
\]
where for $X\subseteq A(a',a)\times B(b,b')$, $\alpha\in A(a'',a')$ and $\beta\in B(b',b'')$, the subset $(\alpha\cdot X\cdot\beta)\subseteq A(a'',a)\times B(b,b'')$ is defined by
\[
(\alpha\cdot X\cdot\beta):=\{(\alpha;\alpha_0,\beta_0;\beta)\mid(\alpha_0,\beta_0)\in X\}.
\]
Note that $\widetilde G$ has a single orbit, i.e., for every $x\in\widetilde G(a,b)$ and $y\in\widetilde G(a',b')$, there exist $\alpha\in A(a',a)$ and $\beta\in B(b,b')$ such that $y=\alpha\cdot x\cdot\beta$. This profunctor cannot be decomposed by $+$ any further.
Every profunctor can be written as a sum of profunctors of the form $\widetilde G$. Given a profunctor $P:A\pf B$, let $(x_i\in P(a_i,b_i))_{i\in I}$ be a family of representatives of the equivalence classes defined by the $(A^{\mathrm{op}}\times B)$-action. That is, for every $x\in P(a,b)$, there exists a unique $i\in I$ such that $x=\alpha\cdot x_i\cdot\beta$ for some $\alpha\in A(a,a_i)$ and $\beta\in B(b_i,b)$. Then $P\cong\coprod_{i\in I}\widetilde{\operatorname{fix}(x_i)}$.
\paragraph{Proposition 14.} A profunctor $P:A\pf B$ can be expressed as
\[
P\cong\coprod_{i\in I}\widetilde G_i,\qquad
\forall i\in I.\ G_i\le A^{\mathrm{op}}(a_i,a_i)\times B(b_i,b_i)
\]
for some family $(G_i)_{i\in I}$. Furthermore,
$(\coprod_i\widetilde G_i);(\coprod_j\widetilde H_j)\cong\coprod_{i,j}(\widetilde G_i;\widetilde H_j)$.
\paragraph{Proposition 15.} For $G\le A^{\mathrm{op}}(a,a)\times B(b,b)$ and $H\le B^{\mathrm{op}}(b',b')\times C(c,c)$, we have $(\widetilde G;\widetilde H)\ne\varnothing$ if and only if $b\cong b'$. Here $\varnothing:A\pf C$ is the profunctor such that $\varnothing(a,c)=\varnothing$.
Note that the composite $(\widetilde G;\widetilde H)$ for $G\le A^{\mathrm{op}}(a,a)\times B(b,b)$ and $H\le B^{\mathrm{op}}(b,b)\times C(c,c)$ is not necessarily of a single orbit. It can be a sum $\coprod_{i\in I}\widetilde K_i$ with $\#I>1$.
\section*{Author source: complete single-orbit comparison (34:10--11)}
\paragraph{Lemma 16.} Let $A,B,C$ be locally finite groupoids equipped with factors $\Phi,\Psi,\Theta$, and $K\le A^{\mathrm{op}}(a,a)\times B(b,b)$ and $H\le B^{\mathrm{op}}(b,b)\times C(c,c)$. Then, for every $a'\in A$ and $c'\in C$,
\[
(\norm{\widetilde K};\norm{\widetilde H})_{[a'],[c']}\le
\norm{\widetilde K;\widetilde H}_{[a'],[c']}.
\]
The equality holds if and only if
\[
\{\beta_1;\beta;\beta_1^{-1}\mid(\mathrm{id},\beta)\in K\}
\cap\{\beta_2;\beta;\beta_2^{-1}\mid(\beta,\mathrm{id})\in H\}
=\{\mathrm{id}\}
\]
for every $\beta_1,\beta_2\in B(b,b)$.
\paragraph{Proof.} Assume $a'\in A$ and $c'\in C$. When $\neg(a'\cong a)$, then $\widetilde K(a',b')=\varnothing$ for every $b'\in B$. So we have $\norm{\widetilde K}_{[a'],[b']}=0$ for every $b'$ and $\coprod_{b'\in B}\widetilde K(a',b')\times\widetilde H(b',c')=\varnothing$. The former implies $(\norm{\widetilde K};\norm{\widetilde H})_{[a'],[c']}=0$, and the latter implies $\norm{\widetilde K;\widetilde H}_{[a'],[c']}=0$. So $\norm{\widetilde K};\norm{\widetilde H}$ and $\norm{\widetilde K;\widetilde H}$ coincide on $([a'],[c'])$-entry when $\neg(a'\cong a)$. The case that $\neg(c'\cong c)$ is similar.
We calculate $(\norm{\widetilde K};\norm{\widetilde H})_{[a],[c]}$. Since $\widetilde K$ is the collection of cosets of $K$,
\[
\#\widetilde K(a,b)=\frac{\#(A^{\mathrm{op}}(a,a)\times B(b,b))}{\#K}.
\]
So
\[
\norm{\widetilde K}_{[a],[b]}=
\frac{\#(A^{\mathrm{op}}(a,a)\times B(b,b))}{\#K\,\Phi^\ominus(a)\Psi^\oplus(b)}
=\frac{\Phi^\oplus(a)\Psi^\ominus(b)}{\#K}.
\]
The above formula, together with $\norm{\widetilde K}_{[a],[b']}=0$ for $\neg(b'\cong b)$, concludes
\[
(\norm{\widetilde K};\norm{\widetilde H})_{[a],[c]}
=\frac{\Phi^\oplus(a)\,\#B(b,b)\,\Theta^\ominus(c)}{\#K\#H}.
\]
We calculate $\norm{\widetilde K;\widetilde H}_{[a],[c]}$. Recall that $(\widetilde K;\widetilde H)(a,c)\cong(\widetilde K(a,b)\times\widetilde H(b,c))/{\sim}$, where $(x,y)\sim(x',y')$ if and only if $(x\cdot\beta,\beta^{-1}\cdot y)=(x',y')$ for some $\beta\in B(b,b)$. For every $x\in\widetilde K(a,b)$ and $y\in\widetilde H(b,c)$, the equivalence class $[(x,y)]_\sim=\{(x\cdot\beta,\beta^{-1}\cdot y)\mid\beta\in B(b,b)\}$ has at most $\#B(b,b)$ elements. Therefore
\[
\begin{aligned}
\#(\widetilde K;\widetilde H)(a,c)&\ge\frac{\#\widetilde K(a,b)\,\#\widetilde H(b,c)}{\#B(b,b)}\\
&=\frac{\#(A^{\mathrm{op}}(a,a)\times B(b,b))}{\#K}
 \times\frac{\#(B^{\mathrm{op}}(b,b)\times C(c,c))}{\#H}
 \times\frac1{\#B(b,b)}.
\end{aligned}
\]
So
\[
\norm{\widetilde K;\widetilde H}_{[a],[c]}
=\frac{\#(\widetilde K;\widetilde H)(a,c)}{\Phi^\ominus(a)\Theta^\oplus(c)}
\ge\frac{\Phi^\oplus(a)\,\#B(b,b)\,\Theta^\ominus(c)}{\#K\#H}.
\]
Assume the additional condition in the statement. Let $x=\alpha\cdot K\cdot\beta_1$ and $y=\beta_2\cdot H\cdot\gamma$ and suppose $(x,y)=(x\cdot\beta,\beta^{-1}\cdot y)$. Then $\alpha\cdot K\cdot\beta_1=\alpha\cdot K\cdot\beta_1\cdot\beta$, so $K=K\cdot\beta_1\cdot\beta\cdot\beta_1^{-1}$. Similarly, $\beta_2\cdot H\cdot\gamma=\beta^{-1}\cdot\beta_2\cdot H\cdot\gamma$ implies $H=\beta_2^{-1}\cdot\beta^{-1}\cdot\beta_2\cdot H$. This means that $(\mathrm{id},\beta_1;\beta;\beta_1^{-1})\in K$ and $(\beta_2^{-1};\beta^{-1};\beta_2,\mathrm{id})\in H$. By the additional assumption, $\beta=\mathrm{id}$. Therefore, $(x,y)=(x\cdot\beta,\beta^{-1}\cdot y)$ implies $\beta=\mathrm{id}$, and thus the equivalence class $[(x,y)]_\sim$ consists of $\#B(b,b)$ elements. So
\[
\#(\widetilde K;\widetilde H)(a,c)=
\frac{\#\widetilde K(a,b)\,\#\widetilde H(b,c)}{\#B(b,b)}.
\]
A straightforward calculation shows that $\norm{\widetilde K;\widetilde H}_{[a],[c]}=(\norm{\widetilde K};\norm{\widetilde H})_{[a],[c]}$.
Conversely, assume that $\{\beta_1;\beta;\beta_1^{-1}\mid(\mathrm{id},\beta)\in K\}\cap\{\beta_2;\beta;\beta_2^{-1}\mid(\beta,\mathrm{id})\in H\}$ contains $\beta\ne\mathrm{id}$ for some $\beta_1,\beta_2\in B(b,b)$. Then $K=K\cdot\beta_1\cdot\beta\cdot\beta_1^{-1}$ and hence $K\cdot\beta_1=K\cdot\beta_1\cdot\beta$. Similarly, since $H=\beta_2\cdot\beta\cdot\beta_2^{-1}\cdot H$, we have $\beta^{-1}\cdot\beta_2^{-1}\cdot H=\beta_2^{-1}\cdot H$. Let $x=K\cdot\beta_1$ and $y=\beta_2^{-1}\cdot H$. Then $(x,y)=(x\cdot\beta,\beta^{-1}\cdot y)$, so the equivalence class $[(x,y)]_\sim$ consists of at most $\#B(b,b)/2$ elements. Recall the inequality
\[
\#(\widetilde K;\widetilde H)(a,c)\ge
\frac{\#\widetilde K(a,b)\,\#\widetilde H(b,c)}{\#B(b,b)}.
\]
Since $\#\widetilde K(a,b)$ and $\#\widetilde H(b,c)$ are finite, the equality holds only if all the equivalence classes $[(x,y)]_\sim$ consist of $\#B(b,b)$ elements. Therefore
\[
\#(\widetilde K;\widetilde H)(a,c)>
\frac{\#\widetilde K(a,b)\,\#\widetilde H(b,c)}{\#B(b,b)}
\]
and hence $\norm{\widetilde K;\widetilde H}_{[a],[c]}>(\norm{\widetilde K};\norm{\widetilde H})_{[a],[c]}$.\hfill$\triangleleft$
\section*{Author source: complete general-case comparison (34:11--12)}
By Proposition 14, the argument for the single-orbit case extends to the general case.
\paragraph{Lemma 17.} Let $P:A\pf B$ and $Q:B\pf C$, and assume factors for $A,B,C$. If $\mathcal C(P(a,-))\perp\mathcal C(Q(-,c))$ for every $a\in A$ and $c\in C$, then $\norm P;\norm Q=\norm{P;Q}$.
\paragraph{Proof.} By Proposition 14, profunctors $P$ and $Q$ can be described as $P\cong\coprod_{i\in I}\widetilde H_i$ and $Q\cong\coprod_{j\in J}\widetilde K_j$ where $H_i\le A^{\mathrm{op}}(a_i,a_i)\times B(b_i,b_i)$ and $K_j\le B^{\mathrm{op}}(b_j,b_j)\times C(c_j,c_j)$. Note that $(P;Q)\cong\coprod_{i,j}(\widetilde H_i;\widetilde K_j)$ and $\norm{P;Q}=\sum_{i,j}\norm{\widetilde H_i;\widetilde K_j}$. So it suffices to show that $\norm{\widetilde H_i;\widetilde K_j}=\norm{\widetilde H_i};\norm{\widetilde K_j}$ for every $i$ and $j$.
We prove that the additional condition in Lemma 16 holds for every pair $(H_i,K_j)$ such that $b_i\cong b_j$. We can assume without loss of generality that $b_i=b_j$. Let $b=b_i=b_j$ and $\beta_1,\beta_2\in B(b,b)$. Assume that $\beta\in B(b,b)$ belongs to $\{\beta_1;\beta;\beta_1^{-1}\mid(\mathrm{id},\beta)\in H_i\}\cap\{\beta_2;\beta;\beta_2^{-1}\mid(\beta,\mathrm{id})\in K_j\}$. Then $H_i\cdot\beta_1$ is a fixed-point of the action $(-)\cdot\beta$, so $\beta\in(\mathcal C(\widetilde H_i(a_i,-)))_b\subseteq(\mathcal C(P(a_i,-)))_b$. Similarly, since $\beta_2^{-1}\cdot K_j$ is a fixed-point of the action $\beta\cdot(-)$, we have $\beta\in(\mathcal C(\widetilde K_j(-,c_j)))_b\subseteq(\mathcal C(Q(-,c_j)))_b$. By the orthogonality assumption $\mathcal C(P(a_i,-))\perp\mathcal C(Q(-,c_j))$, we have $\beta=\mathrm{id}$ as expected.\hfill$\triangleleft$
\section*{Author source: completeness (34:12)}
We have seen that the orthogonality of the symmetries of profunctors provides a sufficient condition for the matrix representation to behave functorially. In fact, this condition is almost necessary: as long as the result of matrix multiplication is finite, the functoriality of the matrix representation entails the orthogonality of the symmetry.
\paragraph{Theorem 20.} Let $P:A\pf B$ and $Q:B\pf C$ be profunctors, and suppose that $\norm P;\norm Q=\norm{P;Q}$. Then, for every $a\in A$ and $c\in C$, we have $\mathcal C(P(a,-))\perp\mathcal C(Q(-,c))$ provided that $\norm{P;Q}_{[a],[c]}<\infty$.
\paragraph{Proof.} By Proposition 14, profunctors $P$ and $Q$ can be described as $P\cong\coprod_{i\in I}\widetilde H_i$ and $Q\cong\coprod_{j\in J}\widetilde K_j$ where $H_i\le A^{\mathrm{op}}(a_i,a_i)\times B(b_i,b_i)$ and $K_j\le B^{\mathrm{op}}(b_j,b_j)\times C(c_j,c_j)$. Note that $(P;Q)\cong\coprod_{i,j}(\widetilde H_i;\widetilde K_j)$ and $\norm{P;Q}=\sum_{i,j}\norm{\widetilde H_i;\widetilde K_j}$.
Assume for contradiction that the creed orthogonality fails for some $a$ and $c$. Then the orthogonality fails for $(\widetilde H_{i_0},\widetilde K_{j_0})$ for some $(i_0,j_0)$. By Lemma 16,
$(\norm{\widetilde H_{i_0}};\norm{\widetilde K_{j_0}})_{[a],[c]}<
\norm{\widetilde H_{i_0};\widetilde K_{j_0}}_{[a],[c]}$,
i.e.,
\[
(\norm{\widetilde H_{i_0}};\norm{\widetilde K_{j_0}})_{[a],[c]}+\epsilon
=\norm{\widetilde H_{i_0};\widetilde K_{j_0}}_{[a],[c]}
\]
for some $\epsilon>0$. Since $(\norm{\widetilde H_i};\norm{\widetilde K_j})_{[a],[c]}\le\norm{\widetilde H_i;\widetilde K_j}_{[a],[c]}$ for every $i$ and $j$, we have
\[
\left(\sum_{i,j}\norm{\widetilde H_i};\norm{\widetilde K_j}\right)_{[a],[c]}+\epsilon
\le\left(\sum_{i,j}\norm{\widetilde H_i;\widetilde K_j}\right)_{[a],[c]}.
\]
Since $\sum_{i,j}\norm{\widetilde H_i;\widetilde K_j}<\infty$, we have
\[
\left(\sum_{i,j}\norm{\widetilde H_i};\norm{\widetilde K_j}\right)_{[a],[c]}
<\left(\sum_{i,j}\norm{\widetilde H_i;\widetilde K_j}\right)_{[a],[c]},
\]
a contradiction.\hfill$\triangleleft$
